\documentclass[a4paper,12pt]{article}
\usepackage[T2A]{fontenc}
\usepackage[utf8]{inputenc}
\usepackage[russian]{babel}
\usepackage[left=23mm,right=23mm,top=24mm,bottom=24mm,headheight=16pt]{geometry}
\usepackage{amsmath,amssymb,booktabs,tabularx,xcolor,tikz,fancyhdr,microtype}
\usepackage[colorlinks=true,linkcolor=teal!60!black]{hyperref}
\definecolor{accent}{HTML}{008080}
\definecolor{darkaccent}{HTML}{004D4D}
\setlength{\parindent}{0pt}
\setlength{\parskip}{3pt}
\pagestyle{fancy}
\fancyhf{}
\fancyhead[L]{\small Правильные пирамиды}
\fancyhead[R]{\small Николь Саркисьянц}
\fancyfoot[C]{\thepage}
\renewcommand{\headrulewidth}{0.35pt}
\newcommand{\sub}[1]{\par\smallskip\textbf{\color{darkaccent}#1}\par}
\newcommand{\err}[1]{\par\textcolor{red!70!black}{\textbf{Ошибка:}} #1\par}
\newcommand{\idea}[1]{\par\textcolor{blue!60!black}{\textbf{Ключевая идея:}} #1\par}
\newcommand{\result}[1]{\par\textcolor{teal!60!black}{\textbf{Итог:}} #1\par}
\newcommand{\task}[2]{\clearpage\phantomsection\label{task:#1}{\Large\color{darkaccent}\bfseries Задание № #1. #2}\par\smallskip{\color{accent}\rule{\linewidth}{0.45pt}}\par}
\begin{document}
\begin{titlepage}
\centering
\vspace*{0.20\textheight}
{\Large\color{darkaccent} МАТЕМАТИКА\par}
\vspace{15mm}
{\LARGE Ревью домашней работы\par}
\vspace{10mm}
{\Large\color{accent} Правильные пирамиды\par}
\vfill
{\large Николь Саркисьянц\par}
\vspace{9mm}25 ноября 2026 г.\par
\vspace{14mm}\small Лёвин Артём Александрович\par эксклюзивно для Николь Саркисьянц\par
\vspace{15mm}
\begin{tikzpicture}\draw[accent,line width=0.7pt] (0,0) .. controls (2,0.8) and (4,-0.8) .. (6,0) .. controls (8,0.8) and (10,-0.8) .. (12,0);\end{tikzpicture}
\end{titlepage}
\setcounter{page}{1}
{\Large\color{darkaccent}\bfseries Сводная таблица}\par\medskip
Ревью выполнено по тренировочному блоку пособия «Правильные пирамиды» от 19 ноября 2026 г.\par\medskip
\renewcommand{\arraystretch}{1.4}
\small
\begin{tabularx}{\linewidth}{@{}c >{\raggedright\arraybackslash}p{26mm} >{\raggedright\arraybackslash}X >{\raggedright\arraybackslash}p{31mm} >{\raggedright\arraybackslash}p{32mm}@{}}
\toprule
№ & Тема & Суть ошибки & Ваш ответ & Правильный ответ\\
\midrule
\hyperref[task:6]{6} & Объём & Формула для диагонали квадратного основания выбрана неверно; решение не закончено. & Не указан & $50$\\
\hyperref[task:9]{9} & Угол между плоскостями & Апофема обозначена буквой высоты; итоговые числа верны. &  $288$;\newline $144+144\sqrt2$ & $288$;\newline $144+144\sqrt2$\\
\hyperref[task:10]{10} & Масштаб & Записано начало первой части; вычислений и второй части нет. & Не указан & $4$; $8$; $12$ раз\\
\bottomrule
\end{tabularx}
\normalsize
\vspace{12mm}
\textbf{Рекомендация.} Различайте высоту пирамиды, апофему и боковое ребро. Перед вычислением подписывайте на чертеже прямоугольный треугольник, из которого следует формула. При изменении размеров отслеживайте, сколько линейных множителей входит в площадь и объём.
\clearpage
\thispagestyle{empty}
\vspace*{\fill}
\begin{center}{\Huge\color{darkaccent} Разбор ошибок}\end{center}
\vspace*{\fill}
\task{6}{Объём пирамиды}
\sub{Условие} Диагональ квадратного основания равна $10$, боковое ребро равно $\sqrt{34}$. Найдите объём правильной пирамиды.
\sub{Ваше решение и ответ} Записано $V=\frac13 S_{\text{осн}}H$, далее использовано соотношение $10^2=3a^2/4$ и получено $S_{\text{осн}}=400/3$. Высота и объём не найдены; ответ отсутствует.
\err{Неподходящая формула для основания; решение не завершено.}
Последний правильный шаг --- формула объёма. Первый неверный шаг --- применение равенства $d^2=3a^2/4$. Такое выражение соответствует квадрату высоты равностороннего треугольника со стороной $a$. Для диагонали квадрата по теореме Пифагора выполняется $d^2=a^2+a^2=2a^2$.
\idea{Используйте диагональ дважды: для площади квадрата и для расстояния от центра основания до вершины.}
\sub{Исправление и полное решение}
Пусть $S$ --- вершина пирамиды, $O$ --- центр квадрата, $A$ --- вершина основания. Из правильности пирамиды следует $SO\perp OA$.
\[
S_{\text{осн}}=a^2=\frac{d^2}{2}=\frac{100}{2}=50,
\qquad OA=\frac d2=5.
\]
В прямоугольном треугольнике $SOA$ гипотенуза $SA=\sqrt{34}$. Высота $H=SO$ является катетом. Поэтому
\[
H=\sqrt{SA^2-OA^2}=\sqrt{34-25}=3.
\]
Подставим найденные площадь и высоту в уже правильно записанную формулу:
\[V=\frac13\cdot50\cdot3=50.\]
\sub{Проверка} $3^2+5^2=34$; высота положительна, пирамида существует. Объём составляет треть объёма призмы с теми же основанием и высотой.
\result{$V=50$.}
\sub{Задача на закрепление} Диагональ основания правильной четырёхугольной пирамиды равна $12$, боковое ребро равно $10$. Найдите объём.
\task{9}{Высота и апофема}
\sub{Условие} Сторона основания правильной четырёхугольной пирамиды равна $12$, угол между боковой гранью и плоскостью основания равен $45^\circ$. Найдите объём и полную площадь поверхности.
\sub{Ваше решение и ответ} Получены $S_{\text{осн}}=144$ и запись $H=\sqrt{6^2+6^2}=6\sqrt2$. Затем в объём подставлено $6$, а в боковую площадь --- $6\sqrt2$. Ответ: $288$; $144+144\sqrt2$.
\err{Высота и апофема смешаны в обозначениях. Численные ответы верны.}
Последний правильный шаг --- площадь основания $12^2=144$. Первая ошибочная запись --- $H=6\sqrt2$: корень из суммы квадратов двух катетов даёт гипотенузу, которой в нужном сечении является апофема. При вычислении объёма Вы использовали верную высоту, поэтому исправлять итоговые числа не требуется.
\idea{Угол между гранью и основанием измеряется в сечении, перпендикулярном ребру основания.}
\begin{center}
\begin{tikzpicture}[scale=0.65]
\draw[thick] (0,0) node[below]{$O$} -- (4,0) node[below right]{$M$} -- (0,4) node[above]{$S$} -- cycle;
\draw (0,0.35)--(0.35,0.35)--(0.35,0);
\draw (3.2,0) arc[start angle=180,end angle=135,radius=0.8];
\node at (2.7,0.4) {$45^\circ$};
\node[left] at (0,2) {$H$};\node[below] at (2,0) {$6$};\node[above right] at (2,2) {$h$};
\end{tikzpicture}
\end{center}
\sub{Исправление и полное решение}
$O$ --- центр основания, $M$ --- середина его стороны. Тогда $OM=6$, $SO\perp OM$, а $\angle SMO=45^\circ$. В прямоугольном равнобедренном треугольнике $SOM$ катеты равны: $SO=OM=6$. Следовательно,
\[
H=6,\qquad h=SM=\sqrt{6^2+6^2}=6\sqrt2.
\]
Периметр основания равен $48$. Отсюда
\[
V=\frac13\cdot144\cdot6=288,
\qquad S_{\text{бок}}=\frac12\cdot48\cdot6\sqrt2=144\sqrt2.
\]
\sub{Проверка и итог} Апофема длиннее высоты: $6\sqrt2>6$; полная площадь включает основание. \result{$V=288$, $S_{\text{полн}}=144+144\sqrt2$.}
\sub{Задача на закрепление} Сторона основания правильной четырёхугольной пирамиды равна $10$, угол между боковой гранью и основанием равен $45^\circ$. Найдите объём и полную площадь поверхности.
\task{10}{Изменение размеров}
\sub{Условие} В правильном тетраэдре все рёбра увеличили в два раза. Во сколько раз изменились площадь полной поверхности и объём? Затем у другой пирамиды сторону квадратного основания увеличили в два раза, а высоту в три раза. Во сколько раз изменился её объём?
\sub{Ваше решение и ответ} Записано $a_2=2a_1$, обозначены отношения площадей и объёмов. Дальнейших вычислений и второй части нет. Ответ отсутствует.
\err{Решение остановилось на записи данных.}
Последний правильный шаг --- обозначение $a_2=2a_1$. Следующий необходимый шаг пропущен: выразить отношения площадей и объёмов через изменение линейных размеров. Арифметической ошибки здесь нет; требуется завершить оба случая.
\idea{При подобии площадь изменяется как квадрат коэффициента подобия, объём --- как его куб. При независимом изменении основания и высоты используйте формулу объёма.}
\sub{Исправление и полное решение}
\textbf{Часть 1.} Каждая грань правильного тетраэдра --- равносторонний треугольник. Если ребро равно $a$, его полная площадь равна четырём площадям таких треугольников:
\[S_{\text{полн}}=4\cdot\frac{a^2\sqrt3}{4}=a^2\sqrt3.\]
Поэтому $S_2/S_1=(2a_1)^2/a_1^2=4$. При подобии высота тетраэдра также удваивается. Площадь основания увеличивается в четыре раза, откуда
\[\frac{V_2}{V_1}=\frac{\frac13(4B_1)(2H_1)}{\frac13 B_1H_1}=8,\]
где $B_1$ в этой дроби обозначает площадь основания первого тетраэдра.

\textbf{Часть 2.} Для пирамиды с квадратным основанием $V_1=a^2H/3$. После изменения размеров $a_2=2a$, $H_2=3H$, поэтому
\[\frac{V_2}{V_1}=\frac{(2a)^2\cdot3H/3}{a^2H/3}=4\cdot3=12.\]
\sub{Проверка} В первом случае три линейных множителя увеличились вдвое; во втором два множителя увеличились вдвое, третий --- втрое.
\result{Тетраэдр: площадь увеличилась в $4$ раза, объём в $8$ раз. Другая пирамида: объём увеличился в $12$ раз.}
\sub{Задача на закрепление} Все рёбра правильного тетраэдра увеличили в три раза. Найдите коэффициенты изменения полной площади и объёма. У другой пирамиды сторону квадратного основания увеличили в три раза, высоту --- в два раза. Во сколько раз увеличился её объём?
\fancyfoot[C]{\footnotesize Лёвин Артём Александрович эксклюзивно для Николь Саркисьянц\quad\thepage}
\end{document}
